% Generated draft: requires mathematical and source review.
\documentclass[11pt]{article}
\usepackage[margin=25mm]{geometry}
\usepackage{fontspec}
\setmainfont{Latin Modern Roman}
\usepackage{amsmath,amssymb,mathtools,unicode-math}
\setmathfont{Latin Modern Math}
\usepackage{enumitem,needspace,xurl,hyperref}
\setlength{\parindent}{0pt}
\setlength{\parskip}{5pt}
\setlength{\emergencystretch}{3em}
\newcommand{\N}{\mathbb{N}}
\newcommand{\Z}{\mathbb{Z}}
\newcommand{\Q}{\mathbb{Q}}
\newcommand{\R}{\mathbb{R}}
\newcommand{\C}{\mathbb{C}}
\newcommand{\F}{\mathbb{F}}
\newcommand{\A}{\mathbb{A}}
\newcommand{\PP}{\mathbb{P}}
\newcommand{\E}{\mathbb{E}}
\newcommand{\eps}{\varepsilon}
\newcommand{\dd}{\,\mathrm{d}}
\begin{document}
\section*{UnsolvedMath ID: 30006076}
\textbf{Source problem code:} catalogue entry
\subsection{Definitions and mathematical statement}
For the source's stated scope, let \(X/\Q\) be a hyperbolic curve and let \(p\) be a prime of good reduction. Under the usual convention, a hyperbolic curve is a smooth geometrically connected curve whose smooth projective compactification has geometric genus \(g\) and whose geometric complement consists of \(r\) points, with \(2g-2+r>0\). The supplied excerpt does not specify whether the rational-point formulation assumes that \(X\) is proper; this is an unresolved hypothesis.

Write \(X(\Q_p)_n\subseteq X(\Q_p)\), for positive integers \(n\), for the level-\(n\) refined Chabauty--Kim sets intended by the source. The supplied record does not give their construction or the auxiliary choices and local conditions defining the refinement. These symbols therefore designate the source's sets rather than a fully specified construction in this draft. Identify \(X(\Q)\) with its image in \(X(\Q_p)\).

The rational-point assertion attributed to Kim, labeled Conjecture 4 in the excerpt, is the following: for every curve \(X\) and prime \(p\) satisfying the intended hypotheses, there exists an integer \(n_0\geq 1\) such that
\[
\forall n\in\Z,\qquad n\geq n_0\ \Longrightarrow\ X(\Q_p)_n=X(\Q).
\]
The threshold may depend on \(X\), \(p\), and the data defining the refined sets; no uniform threshold is asserted. This is the supplied conjectural formulation, subject to the unresolved hypotheses and definition above.

The excerpt separately labels as Conjecture 3 the assertion
\[
\exists n_1\geq 1\ \forall n\geq n_1,\qquad \#X(\Q_p)_n<\infty.
\]
It states that this finiteness follows from the Bloch--Kato conjecture, citing its reference [4], and attributes eventual equality to [7, Conjecture 3.1]. Those references and the precise Bloch--Kato hypothesis are not supplied. The finiteness and equality assertions are recorded separately. In particular, the condition of being hyperbolic alone allows affine curves with infinitely many rational points, so the missing scope cannot be resolved by silently imposing both assertions on every affine hyperbolic curve.

The excerpt also gives an \(S\)-integral variant for affine hyperbolic curves. Under the conventional interpretation that \(S\) is a finite set of rational primes, put
\[
\Z_S=\Z[\ell^{-1}:\ell\in S].
\]
An integral model \(\mathcal X/\Z_S\) with generic fiber \(X\) is needed to interpret the source's notation \(X(\Z_S)\) and \(X(\Z_p)\); no such model is specified. For a prime \(p\notin S\) of good reduction, the source describes refined sets satisfying
\[
\mathcal X(\Z_S)\subseteq X(\Z_p)_{S,n}\subseteq\mathcal X(\Z_p)
\]
and states the conjectural equality
\[
\exists n_0\geq 1\ \forall n\geq n_0,\qquad X(\Z_p)_{S,n}=\mathcal X(\Z_S).
\]
Here the integral points are viewed through their natural map to \(\mathcal X(\Z_p)\). This variant remains dependent on the unspecified model and refined local conditions.
\subsection{Short English statement}
Kim's conjecture, as supplied, asks whether the refined Chabauty--Kim sets equal the rational points at every sufficiently large level, with a threshold depending on the curve and prime. The excerpt also states a separate finiteness conjecture and an affine \(S\)-integral analogue. The rational-point scope, the construction of the refined sets, and the integral model require clarification before this can be a fully self-contained statement.
\subsection{Sources}
\begin{itemize}
\item Local UnsolvedMath registry, supplied complete record: permanent ID \(30006076\), exact source problem code \texttt{catalogue entry}, title ``Kim's Nonabelian Chabauty Conjecture.'' The record lists proposed year 2024 and no proposer; Kim is named in the conjecture's attribution.
\item Supplied source citation: \emph{Anabelian Geometry and Representations of Fundamental Groups} (2025), Oberwolfach Reports, pp. 2621--2670, DOI \texttt{10.4171/owr/2024/45}. \url{https://doi.org/10.4171/owr/2024/45}. The extracted passage includes a page marker for p. 2647 and labels the principal assertion Conjecture 4.
\item Supplied literature source: ``Oberwolfach Report 45/2024.'' \url{https://doi.org/10.4171/owr/2024/45}.
\item Supplied literature source: Betts--Kumpitsch--Litt, \emph{Chabauty--Kim and the Section Conjecture}. \url{https://arxiv.org/abs/2305.09462}.
\item Literature assessment dated exactly \texttt{2026-08-22}: the supplied curation labels the problem open and says that eventual equality remains a general conjecture, with finiteness and equality proved only in restricted families and levels. This is supplied dated curation, not an independently verified literature assessment.
\item The record labels the statement \texttt{corrected\_verified} and reports that the official Oberwolfach report was checked, the question isolated, and notation repaired. These are supplied verification claims; no independent source consultation is claimed here.
\item The excerpt cites [4] for the Bloch--Kato implication, [7, Conjecture 3.1] for eventual equality, and [2, Corollary 3.13] for a surrounding claim about real places. Their bibliographic identities are absent from the supplied record and have not been inferred from the separately listed literature sources.
\end{itemize}
\end{document}
